\section{10. Estado actual del proyecto}
\small
\begin{longtable}{|L{0.30\textwidth}|L{0.20\textwidth}|L{0.43\textwidth}|}\hline
\rowcolor{Navy}\THead{Componente} & \THead{Estado actual} & \THead{Resultado} \\ \hline\endfirsthead
\rowcolor{Navy}\THead{Componente} & \THead{Estado actual} & \THead{Resultado} \\ \hline\endhead
Servidor multimedia & Operativo & Recepción y distribución de contenido en funcionamiento. \\ \hline
ENLACE y EJTV & Señales principales & Casos centrales de operación. \\ \hline
Cable Centro & Primer Reader real & Distribución hacia receptor externo. \\ \hline
MediaMTX & Operativo & Motor actual de ingesta y distribución. \\ \hline
SRT / RTMP / RTSP / HLS / WebRTC & Disponibles & Diversidad de mecanismos de transporte y entrega. \\ \hline
NOC Terminal & Funcional & Supervisión técnica de infraestructura, señales, sesiones, clientes, Health, eventos y alarmas. \\ \hline
NOC Web & En evolución & Interfaz web destinada a ampliar el acceso visual para operadores, administradores y usuarios autorizados. \\ \hline
Identificación de conexiones & Implementada & Permite conocer Readers/sesiones, señal, protocolo, origen y actividad según la evidencia disponible. \\ \hline
Health por protocolo & Implementado & Especialización para cinco protocolos principales. \\ \hline
Análisis audiovisual avanzado & Previsto & Detección futura de presencia de audio/video, congelamiento, pantalla negra y otras anomalías de contenido. \\ \hline
Multiviewer & Previsto & Monitoreo visual simultáneo y en tiempo real de las señales entrantes. \\ \hline
Eventos / alarmas / historia & Implementados & Trazabilidad y soporte operacional. \\ \hline
Validación física & En curso & Pruebas sobre señales y conexiones reales. \\ \hline
IMPACT CBR A/B & Pendiente externo & Baseline cerrado; requiere coordinación con cliente. \\ \hline
Transcodificación dedicada & Etapa posterior & Capacidad prevista para infraestructura separada. \\ \hline
Escalamiento multi-nodo & Etapa futura & Evolución prevista de la arquitectura. \\ \hline
\end{longtable}
\normalsize
El desarrollo se ejecuta bajo una metodología de inspección, definición de contrato, pruebas, cambios controlados, regresión, validación física y control de versiones. Este enfoque busca que las capacidades incorporadas respondan a evidencia operacional verificable.
